% !PRINTED 494
\weberpage{494}

Ist nämlich
% [sic] Im Druck fehlt in der unteren Zeile von (1) das Komma zwischen $a_1$ und $a_2$; quellentreu beibehalten.
\weq{1}{\pi = \begin{pmatrix} 0, & 1, & 2 \dots m-1 \\ a_0, & a_1 & a_2 \dots a_{m-1} \end{pmatrix},}
so fange man mit einer beliebigen Ziffer, etwa mit 0 an und setze die Reihe
\weq{2}{0, \quad a_0 = b, \quad a_b = c \dots}
so lange fort, bis man auf eine Ziffer zum zweiten Male stösst. Die zuerst wiederkehrende Ziffer muss 0 sein, da zu jeder Ziffer die in der Reihe (2) vorangehende durch (1) eindeutig bestimmt ist. Dann bilden die Ziffern
\weq{3}{(0, b, c \dots)}
einen ersten Cyklus. Sind dadurch noch nicht alle $m$ Ziffern von (1) erschöpft, so greift man aus den übrigen eine heraus und verfährt ebenso, bis alle $m$ Ziffern von (1) in den Cyklen untergebracht sind. Da in jedem solchen Cyklus zu jeder Ziffer, die vorangehende sowohl als die nachfolgende, durch (1) völlig bestimmt ist, so sind auch die Cyklen selbst eindeutig bestimmt. Bei der \gesperrt{Bezeichnung} von $\pi$ durch die Cyklen können aber nicht nur die verschiedenen Cyklen beliebig angeordnet, sondern auch in jedem Cyklus mit einer beliebigen seiner Ziffern angefangen werden. Eine Ziffer, die nicht geändert wird, bildet für sich einen eingliedrigen Cyklus.

Wir wählen ein ganz beliebiges Beispiel, wodurch das Verfahren sofort klar wird:
\[
\pi_1 = \begin{pmatrix} 0, & 1, & 2, & 3, & 4, & 5, & 6, & 7 \\ 2, & 3, & 6, & 7, & 1, & 4, & 5, & 0 \end{pmatrix} = (0, 2, 6, 5, 4, 1, 3, 7)
\]
\[
\pi_2 = \begin{pmatrix} 0, & 1, & 2, & 3, & 4, & 5, & 6, & 7 \\ 4, & 2, & 1, & 3, & 5, & 7, & 0, & 6 \end{pmatrix} = (0, 4, 5, 7, 6)(1, 2)(3)
\]
\[
\pi_3 = \begin{pmatrix} 0, & 1, & 2, & 3, & 4, & 5, & 6, & 7 \\ 7, & 2, & 3, & 6, & 5, & 0, & 1, & 4 \end{pmatrix} = (0, 7, 4, 5)(1, 2, 3, 6).
\]
Das Nebeneinandersetzen der Cyklen ist mit einer Composition in dem bisherigen Sinne gleichbedeutend, nur ist zu bemerken, dass Permutationen, die gar keine gemeinschaftliche Ziffer enthalten, bei der Composition immer vertauscht werden können.

Eingliedrige Cyklen, die nichts ändern, werden in der Bezeichnung auch oft weggelassen.
